\documentclass[notheorems,hyperref={bookmarks=true}]{beamer}

% --- PACKAGES ---
\usepackage[framesassubsections]{beamerprosper}
\usepackage{multirow}
\usepackage[utf8]{inputenc} 
\usepackage[utf8]{vietnam} 
\usepackage{amsmath, amsthm, amssymb, tikz}
\usetikzlibrary{calc, shapes.geometric, positioning, shadows}
\usepackage{graphicx}
\usepackage{xcolor}
\usepackage{booktabs} 
\usepackage[table]{xcolor}

% --- THEME ---
\usetheme{CambridgeUS}
\usecolortheme{rose} 
\usefonttheme[onlylarge]{structurebold}
\usefonttheme{professionalfonts}

\setbeamercolor{title}{fg=red!80!black,bg=blue!20!white}
\setbeamertemplate{navigation symbols}{}
\setbeamertemplate{blocks}[rounded][shadow=true]

\logo{\includegraphics[height=1cm]{logoHust.png}}

\setbeamertemplate{footline}{
  \leavevmode%
  \hbox{%
    \begin{beamercolorbox}[wd=.33\paperwidth,ht=2.5ex,dp=1.5ex,center]{author in head/foot}%
      Nguyễn Tuấn Tài
    \end{beamercolorbox}%
    \begin{beamercolorbox}[wd=.34\paperwidth,ht=2.5ex,dp=1.5ex,center]{title in head/foot}%
      \footnotesize{AI Agent hỗ trợ học ngoại ngữ}
    \end{beamercolorbox}%
    \begin{beamercolorbox}[wd=.33\paperwidth,ht=2.5ex,dp=1.5ex,right]{date in head/foot}%
      \insertframenumber/\inserttotalframenumber \hspace*{2ex} 
    \end{beamercolorbox}%
  }%
  \vskip0pt%
}

\title[AI Agent dạy ngoại ngữ]{
    \Large\bfseries Xây dựng AI Agent hỗ trợ học ngoại ngữ
}

\institute[\scriptsize Khoa Toán - Tin]{
    \small
    \begin{tabular}{l}
        \textbf{Giảng viên hướng dẫn: \textcolor{red!80!black}{TS. Tạ Anh Sơn}} \\
        \textbf{Sinh viên thực hiện: \textcolor{red!80!black}{Nguyễn Tuấn Tài (20216959)}} \\
    \end{tabular} \\[0.5em]
    \textbf{Khoa Toán - Tin, Đại học Bách khoa Hà Nội} \\[0.3em]
    \includegraphics[height=1.5cm]{fami.png}
}

\date[Hà Nội, 06/2026]{\scriptsize\bfseries\color{red!80!black}Hà Nội, 06/2026}

\begin{document}

\begin{frame}
    \maketitle
\end{frame}

\footnotesize

% ==========================================
% SLIDE 1: ĐẶT VẤN ĐỀ
% ==========================================
\begin{frame}{Vấn đề cần giải quyết}
\begin{columns}[c]
    \begin{column}{0.5\textwidth}
        \begin{alertblock}{3 Hạn chế học truyền thống}
            \begin{enumerate}
                \item Chi phí cao (200k-500k/giờ)
                \item Không cá nhân hóa
                \item Bài tập hạn chế
            \end{enumerate}
        \end{alertblock}
        
        \vspace{0.5cm}
        
        \begin{exampleblock}{Giải pháp}
            \textbf{Hệ thống Multi-Agent}
            \begin{itemize}
                \item Miễn phí, 24/7
                \item Phân tích lỗi tức thì
                \item Sinh bài tập vô hạn
            \end{itemize}
        \end{exampleblock}
    \end{column}
    \begin{column}{0.5\textwidth}
        \centering
        \begin{figure}
        \resizebox{0.9\textwidth}{!}{
        \begin{tikzpicture}[
            vertex/.style={circle, fill=red!10, draw=red!70, thick, minimum size=15mm, align=center, font=\tiny\bfseries},
            target/.style={circle, draw=blue!70, fill=blue!10, very thick, dashed, minimum size=22mm, align=center, font=\small\bfseries\color{blue!90}}
        ]
            \node[vertex] (Tutor) at (90:2.8) {Gia sư\\truyền thống};
            \node[vertex] (App) at (210:2.8) {Ứng dụng\\tự học};
            \node[vertex] (Center) at (330:2.8) {Trung tâm\\ngoại ngữ};
            \node[target] (Goal) at (0,0) {GIA SƯ AI\\\normalfont\tiny (Mục tiêu)};
            
            \draw[->, blue!60, thick] (Tutor) -- (Goal);
            \draw[->, blue!60, thick] (App) -- (Goal);
            \draw[->, blue!60, thick] (Center) -- (Goal);
        \end{tikzpicture}
        }
        \end{figure}
    \end{column}
\end{columns}
\end{frame}

% ==========================================
% SLIDE 2: KIẾN TRÚC MULTI-AGENT
% ==========================================
\begin{frame}{Kiến trúc Multi-Agent}
\begin{figure}[H]
\centering
\begin{tikzpicture}[
    node distance=1.3cm,
    agent_box/.style={rectangle, draw=gray!60, fill=blue!8, text width=3.5cm, align=center, minimum height=2cm, rounded corners=5pt, font=\scriptsize},
    arrow/.style={->, >=stealth, thick, color=gray!70}
]

\node[agent_box] (ErrorAgent) at (0, 0) {
    \textbf{\color{red!80}Error Analyzer}\\[0.2cm]
    {\tiny • Nhận diện lỗi\\• Giải thích chi tiết}
};

\node[agent_box, right=of ErrorAgent, fill=green!8] (ExerciseAgent) {
    \textbf{\color{green!80}Exercise Generator}\\[0.2cm]
    {\tiny • Sinh bài tập tự động\\• Thích ứng độ khó}
};

\node[agent_box, right=of ExerciseAgent, fill=purple!8] (WritingAgent) {
    \textbf{\color{purple!80}Writing Evaluator}\\[0.2cm]
    {\tiny • Chấm 4 tiêu chí\\• Phản hồi chi tiết}
};

\node[
    rectangle, draw=orange!60, fill=orange!15, text width=2.5cm, align=center,
    minimum height=1.5cm, rounded corners=5pt, font=\small,
    below=2cm of ExerciseAgent
] (AITutor) {
    \textbf{\color{orange!80}AI Tutor}\\
    \tiny (Điều phối)
};

\draw[arrow] (ErrorAgent.south) -- (AITutor.north west);
\draw[arrow] (ExerciseAgent.south) -- (AITutor.north);
\draw[arrow] (WritingAgent.south) -- (AITutor.north east);

\end{tikzpicture}
\caption{3 Agents chuyên biệt + AI Tutor điều phối}
\end{figure}
\end{frame}

% ==========================================
% KẾT QUẢ 1: MÔ HÌNH TRI THỨC
% ==========================================
\begin{frame}{Kết quả 1: Mô hình tri thức toàn diện}
\begin{columns}[c]
    \begin{column}{0.45\textwidth}
        \begin{block}{190 Chủ đề CEFR}
            \textbf{Phạm vi:} A1 → C2\\[0.3cm]
            \textbf{Mỗi chủ đề:} 5 bài học
            \begin{enumerate}
                \item Grammar
                \item Vocabulary
                \item Practice
                \item Writing
                \item Quiz
            \end{enumerate}
        \end{block}
        
        \begin{alertblock}{Kết quả}
            \textbf{950 bài học} được chuẩn hóa
        \end{alertblock}
    \end{column}
    \begin{column}{0.55\textwidth}
        \centering
        \includegraphics[width=\textwidth]{ũiuhsx.png}
        
        {\scriptsize Giao diện chủ đề bài học}
    \end{column}
\end{columns}
\end{frame}

% ==========================================
% KẾT QUẢ 2: PHÂN TÍCH LỖI
% ==========================================
\begin{frame}{Kết quả 2: Phân tích lỗi thông minh}
\begin{columns}[c]
    \begin{column}{0.5\textwidth}
        \begin{exampleblock}{Error Analyzer}
            \begin{itemize}
                \item \textbf{Phát hiện:} Ngữ pháp/Từ vựng
                \item \textbf{Giải thích:} Tiếng Việt chi tiết
                \item \textbf{Lưu trữ:} Ghi nhớ lỗi sai
            \end{itemize}
        \end{exampleblock}
        
        \vspace{0.5cm}
        
        \begin{block}{Workflow}
            \begin{tikzpicture}[node distance=1cm, box/.style={rectangle, draw, fill=blue!10, text width=2.5cm, align=center, minimum height=0.6cm, rounded corners, font=\tiny}, arrow/.style={->, >=stealth}]
                \node[box] (1) {Người học làm sai};
                \node[box, below=of 1] (2) {Phân tích lỗi};
                \node[box, below=of 2] (3) {Giải thích + Đề xuất};
                
                \draw[arrow] (1) -- (2);
                \draw[arrow] (2) -- (3);
            \end{tikzpicture}
        \end{block}
    \end{column}
    \begin{column}{0.5\textwidth}
        \centering
        \fbox{\includegraphics[width=0.9\textwidth]{jhsavs.png}}
        
        {\scriptsize AI Tutor phân tích lỗi và giải thích}
    \end{column}
\end{columns}
\end{frame}

% ==========================================
% KẾT QUẢ 3: SINH BÀI TẬP
% ==========================================
\begin{frame}{Kết quả 3: Sinh bài tập cá nhân hóa}
\begin{block}{Exercise Generator - Luyện tập vô hạn}
    \begin{columns}[c]
        \begin{column}{0.5\textwidth}
            \textbf{Input:}
            \begin{itemize}
                \item Kỹ năng yếu
                \item Chủ đề + trình độ
            \end{itemize}
            
            \vspace{0.3cm}
            
            \textbf{Output:}
            \begin{itemize}
                \item 5-10 bài tập mới
                \item Độ khó thích ứng
                \item Tập trung đúng lỗi
            \end{itemize}
        \end{column}
        \begin{column}{0.5\textwidth}
            \centering
            \begin{tikzpicture}[
                box/.style={rectangle, draw, fill=green!10, text width=3cm, align=center, minimum height=0.7cm, rounded corners, font=\small},
                arrow/.style={->, >=stealth, thick}
            ]
                \node[box] (1) at (0, 1.5) {\tiny Lỗi sai};
                \node[box] (2) at (0, 0.5) {\tiny Agent sinh bài tập};
                \node[box] (3) at (0, -0.5) {\tiny 5-10 câu mới};
                \node[box] (4) at (0, -1.5) {\tiny Học sinh luyện tập};
                
                \draw[arrow] (1) -- (2);
                \draw[arrow] (2) -- (3);
                \draw[arrow] (3) -- (4);
                \draw[arrow, dashed] (4) to[bend right=60] (1);
            \end{tikzpicture}
        \end{column}
    \end{columns}
\end{block}

\vspace{0.3cm}

\begin{alertblock}{Lợi ích}
    \textbf{Không giới hạn} nguồn luyện tập • Củng cố \textbf{ngay tại chỗ} • Không lặp lại
\end{alertblock}
\end{frame}

% ==========================================
% KẾT QUẢ 4: ĐÁNH GIÁ BÀI VIẾT
% ==========================================
\begin{frame}{Kết quả 4: Đánh giá bài viết tự động}
\begin{columns}[c]
    \begin{column}{0.5\textwidth}
        \begin{exampleblock}{Writing Evaluator}
            \textbf{4 Tiêu chí CEFR:}
            \begin{enumerate}
                \item Ngữ pháp (Grammar)
                \item Từ vựng (Vocabulary)
                \item Nội dung (Content)
                \item Cấu trúc (Structure)
            \end{enumerate}
        \end{exampleblock}
        
        \vspace{0.5cm}
        
        \begin{block}{Phản hồi chi tiết}
            • Chỉ ra từng câu cần sửa\\
            • Gợi ý cách hành văn\\
            • Lưu lịch sử tiến bộ
        \end{block}
    \end{column}
    \begin{column}{0.5\textwidth}
        \centering
        \begin{table}
            \tiny
            \begin{tabular}{|l|c|}
                \hline
                \rowcolor{purple!20} \textbf{Tiêu chí} & \textbf{Điểm} \\
                \hline
                Grammar & 8.0/10 \\
                \hline
                Vocabulary & 7.5/10 \\
                \hline
                Content & 8.5/10 \\
                \hline
                Structure & 8.0/10 \\
                \hline
                \rowcolor{green!20} \textbf{Tổng} & \textbf{8.0/10} \\
                \hline
            \end{tabular}
        \end{table}
        
        \vspace{0.5cm}
        
        {\scriptsize Ví dụ kết quả đánh giá tự động}
    \end{column}
\end{columns}
\end{frame}

% ==========================================
% KẾT QUẢ 5: BỘ NHỚ HỌC TẬP
% ==========================================
\begin{frame}{Kết quả 5: Quản lý bộ nhớ học tập}
\begin{block}{Hệ thống bộ nhớ 2 tầng}
    \begin{columns}[c]
        \begin{column}{0.5\textwidth}
            \begin{exampleblock}{Short-term}
                • Mạch hội thoại\\
                • Ngữ cảnh phiên hiện tại\\
                • Lỗi vừa mắc
            \end{exampleblock}
        \end{column}
        \begin{column}{0.5\textwidth}
            \begin{exampleblock}{Long-term}
                • Hồ sơ lỗi sai\\
                • Xu hướng học tập\\
                • Dashboard thống kê
            \end{exampleblock}
        \end{column}
    \end{columns}
\end{block}

\vspace{0.5cm}

\begin{figure}
\centering
\begin{tikzpicture}[
    box/.style={rectangle, draw, fill=yellow!10, text width=3cm, align=center, minimum height=1cm, rounded corners, font=\small},
    db/.style={cylinder, draw, fill=blue!10, shape border rotate=90, aspect=0.3, minimum height=1cm, minimum width=2cm, font=\small}
]
    \node[box] (user) at (0, 0) {Người học\\tương tác};
    \node[db, right=2cm of user] (stm) {Short-term\\Memory};
    \node[db, right=2cm of stm] (ltm) {Long-term\\Memory\\(PostgreSQL)};
    
    \draw[->, >=stealth, thick] (user) -- (stm);
    \draw[->, >=stealth, thick] (stm) -- (ltm);
    \draw[->, >=stealth, thick, dashed] (ltm) to[bend right=30] node[above, font=\tiny]{Cá nhân hóa} (user);
\end{tikzpicture}
\end{figure}
\end{frame}

% ==========================================
% KẾT QUẢ 6: TRIỂN KHAI PRODUCTION
% ==========================================
\begin{frame}{Kết quả 6: Triển khai thực tế}
\begin{table}
    \centering
    \small
    \begin{tabular}{|l|l|}
        \hline
        \rowcolor{blue!20} \textbf{Thành phần} & \textbf{Công nghệ} \\
        \hline
        Backend API & FastAPI + LangChain \\
        \hline
        Frontend & Streamlit \\
        \hline
        Database & PostgreSQL (Neon Tech) \\
        \hline
        LLM Service & Groq API (Llama 3.1 70B) \\
        \hline
        Hosting & Render.com \\
        \hline
    \end{tabular}
\end{table}

\vspace{0.5cm}

\begin{exampleblock}{Hệ thống đang Live}
    \centering
    \textbf{URL:} \texttt{https://lang-prj.onrender.com}
\end{exampleblock}

\begin{block}{Tính năng hoạt động}
    ✅ 190 chủ đề • ✅ Chat với AI • ✅ Phân tích lỗi • ✅ Sinh bài tập • ✅ Đánh giá bài viết
\end{block}
\end{frame}

% ==========================================
% TỔNG KẾT
% ==========================================
\begin{frame}{Tổng kết}
\begin{exampleblock}{Giải quyết 3 vấn đề đặt ra}
    ✅ \textbf{Chi phí/Thời gian:} Miễn phí, 24/7\\
    ✅ \textbf{Cá nhân hóa:} Phân tích lỗi + bộ nhớ dài hạn\\
    ✅ \textbf{Bài tập hạn chế:} Sinh vô hạn, thích ứng trình độ
\end{exampleblock}

\vspace{0.3cm}

\begin{block}{6 Kết quả kỹ thuật}
    \begin{enumerate}
        \item Kiến trúc Multi-Agent (3 agents + AI Tutor)
        \item Mô hình tri thức (190 chủ đề × 5 bài = 950 bài học)
        \item Error Analyzer (phân tích lỗi 2 cấp độ)
        \item Exercise Generator (sinh bài tập tự động)
        \item Writing Evaluator (chấm 4 tiêu chí CEFR)
        \item Triển khai production (live trên Cloud)
    \end{enumerate}
\end{block}
\end{frame}

% ==========================================
% CẢM ƠN
% ==========================================
\begin{frame}
    \begin{center}
        \LARGE\bfseries Cảm ơn thầy cô và các bạn\\đã lắng nghe!
        
        \vspace{1cm}
        
        \normalsize
        \textbf{Sinh viên thực hiện:} Nguyễn Tuấn Tài\\
        \textbf{MSSV:} 20216959\\
        \textbf{Giảng viên hướng dẫn:} TS. Tạ Anh Sơn
        
        \vspace{0.5cm}
        
        \includegraphics[width=0.3\linewidth]{sloth.jpg}
    \end{center}
\end{frame}

\end{document}
